% !TeX program = xelatex
\documentclass[11pt,a4paper]{article}
\usepackage[margin=23mm]{geometry}
\usepackage{fontspec}
\setmainfont{Latin Modern Roman}
\usepackage{amsmath,amssymb,mathtools}
\usepackage{xcolor,enumitem,booktabs,longtable,array,xurl,hyperref}
\definecolor{muted}{HTML}{53636C}
\hypersetup{colorlinks=true,urlcolor=blue,linkcolor=blue}
\setlength{\parindent}{0pt}
\setlength{\parskip}{5pt}
\newcommand{\R}{\mathbb R}
\newcommand{\E}{\mathbb E}
\newcommand{\N}{\mathbb N}
\newcommand{\ind}{\mathbf 1}
\newcommand{\Var}{\operatorname{Var}}
\newcommand{\Cov}{\operatorname{Cov}}
\newcommand{\supp}{\operatorname{supp}}
\DeclareMathOperator*{\argmin}{arg\,min}
\DeclareMathOperator*{\argmax}{arg\,max}
\newcommand{\entrymeta}[3]{{\sffamily\small\color{muted}\textbf{\texttt{#1}}\quad|\quad #2\par #3\par}}
\newcommand{\Problemref}[1]{\texttt{#1}}
\newcommand{\JELref}[2]{\texttt{#2} (JEL \texttt{#1})}
\begin{document}
\section*{Algorithmic pricing and collusion}
\textbf{Problem ID:} \texttt{L41-109}\quad \textbf{JEL:} \texttt{L41}\par
% BEGIN ECONOMICS STATEMENT
\label{problem:OP-0074}
\label{atlas:prob:I4}

\noindent\textbf{Original identifier:} I4 (atlas item 74). \textbf{Classification:} Algorithmic competition robustness research program. \textbf{Original tractability label:} Medium (editorial, not a complexity result).

\subsection*{Definitions and assumptions}

Consider a specified repeated differentiated-product price game with firms $i=1,\ldots,n$, demand $q_i(p,s)$, marginal costs $c_i$, and common state $s$. Each independently installed algorithm maps its own information history and internal learning state into a price distribution. Define a benchmark $p^N(s)$ as a selected one-shot Nash equilibrium of the same demand and cost game; competition need not mean zero markups. Define average excess profit or price relative to this benchmark, weighted by prespecified quantities. Specify information sharing, vendor overlap, exploration, replacement, entry, and discounting. Independence of installation does not imply statistical independence of observed prices.

This is a research program over declared model families, not a single fully specified game. Each instance $m$ must supply price domains, demand and payoff maps, state transitions, the timing of prices and observations, information histories, algorithm updates and feedback, initialization, and all sources of randomness, together inducing a probability law $\Pr_m$. With entry, declare the population and active-firm indexing, the entry and exit rules, and how the benchmark changes with market composition; the fixed indexing above applies directly to fixed-population instances. A common state need not be observed by every firm. State whether the benchmark uses the actual one-shot information structure or is a hypothetical full-information comparison. Require existence and a declared measurable equilibrium selection on the configurations under consideration. Declare whether $p^N$ is a pure price vector or denotes a mixed equilibrium distribution; in the latter case evaluate benchmark profits, prices and quantities under that distribution rather than substituting mean prices into nonlinear payoff maps.

For a profit-gap claim, declare the finite evaluation index set $I_t$, per-firm profit statistic $\pi_{it}$, benchmark statistic $b_{it}$, and nonnegative weights $w_{it}$, and set
\[
G_t=\sum_{i\in I_t}w_{it}(\pi_{it}-b_{it}).
\]
Specify the weight rule and any normalization, the units of $G_t$, realized versus conditionally expected profits, the conditioning information, and the treatment of shocks, adoption and replacement costs, entrants and inactive firms. Define how benchmark demand and quantities are evaluated and how benchmark and actual states or shocks are coupled. For a price-gap claim, use a separately declared statistic with prices and their benchmark counterparts in place of profits; price and profit gaps are not interchangeable. The displayed aggregate criterion does not require every firm to gain. A firmwise claim must state its separate events and whether their probability guarantee is joint. The tolerance $\varepsilon$ has the units of the chosen gap statistic.

\subsection*{Formal research question}

Characterize algorithm and market classes for which supracompetitive behavior is robust. For each class-level claim, declare a candidate family $\mathcal C$ of complete instances independently of the desired conclusion, including the permitted market primitives, algorithm profiles, initializations, vendor heterogeneity, shocks, turnover and entry processes. Describe the admissible perturbations and their scope; if robustness is local in a metric or topology, supply that structure and the neighborhood. For fixed declared $\varepsilon>0$ and $\alpha\in(0,1)$, seek substantive structural conditions under which
\[
\Pr_m\!\left(\liminf_{T\to\infty}\frac1T\sum_{t=1}^TG_t\geq\varepsilon\right)\geq1-\alpha
\qquad\text{for every }m\in\mathcal C.
\]
Here the same tolerances apply throughout the declared class. Pointwise claims or tolerances depending on primitives must state that dependence and their quantifier order explicitly; they do not establish a common class-wide margin. The probability concerns the declared process randomness, not the coverage error of an empirical confidence procedure. Characterization requires informative restrictions on the inputs, rather than merely naming the instances for which the event holds.

A finite-horizon claim must separately declare an integer burn-in $B\geq0$, an integer evaluation length $H\geq1$, and its horizon quantifiers, for example a bound on
\[
\Pr_m\!\left(\frac1H\sum_{t=B+1}^{B+H}G_t\geq\varepsilon\right).
\]
State whether the claim holds at one specified horizon, at every horizon in a declared set, or simultaneously over that set, and declare any horizon-dependent tolerances. Lack of price convergence does not by itself rule out the asymptotic time-average criterion.

Evaluate robustness to heterogeneous vendors and profitable unilateral deviations separately from the gap event. Each deviation assessment must declare the deviation time and information available then, the admissible replacement algorithms or strategies, replacement costs, rivals' continuing update rules, and the deviator's payoff objective. Define profitability by a positive increase in that declared objective relative to continuing the installed algorithm, or state an explicit approximation tolerance. State whether the assessment is ex ante, conditional on specified histories, or required at every history in a declared domain. Discounted deviation payoffs and undiscounted average gaps may serve different purposes, but their roles must be explicit. Persistence of a positive gap alone is neither an equilibrium nor a no-profitable-deviation guarantee.

Separately identify causal effects of adoption on market outcomes using randomized adoption or credible exogenous variation. Each empirical application must declare the population and market boundaries, adoption versions and timing, the intervention vector or exposure mapping allowing simultaneous adoption and spillovers, comparison regimes, outcome and horizon, target causal estimand, and observable data. State design-specific assignment and identification assumptions, including the assumptions needed for any instrument-specific effect; a design label alone is insufficient. Distinguish identification from estimation and inference, and do not equate an individual adoption effect with the effect of market-wide adoption under interference. No particular intervention, exposure restriction or empirical design is prescribed for the whole program.

\subsection*{Interpretation and scope}

The bare existence question has already received affirmative constructive evidence in particular simulations. The substantive extension concerns when that behavior survives realistic costs, demand shocks, algorithm turnover, and entry. Parallel prices can reflect common costs or competitive response and therefore do not identify tacit cooperation. Nor does an above-Nash average price prove legal collusion: the economic benchmark and legal agreement standard are distinct. Distinguish convergence of individual prices, convergence in distribution, stationarity, and sample-path time averages. A stationary distribution alone establishes neither attraction from the declared initializations nor the displayed time-average probability bound; supply convergence or ergodic arguments appropriate to the claim actually made. Finite-run simulation estimates should report sensitivity to random seeds, initialization, and exploration schedules and cannot prove the displayed probability bound for all environments. Deviation tests need to allow firms to modify algorithms, not merely change one price while retaining a fixed learning rule. Empirical adoption studies additionally confront endogenous timing and local market exposure. A model-based robustness result and a causal market estimate are complementary deliverables with different evidential requirements.

\subsection*{Source audit and status}

[S1] provides Q-learning simulation evidence; [S2], published in 2024, provides German gasoline-market adoption evidence; [S3] analyzes legal-economic issues involving autonomous agents. These source roles motivate a robustness and identification research program beyond the bare existence question; they do not establish that every branch of the program remains unresolved in the current literature. The displayed asymptotic criterion is an editorial target and is not asserted as a theorem or universal empirical fact established by these papers.

\subsection*{References}

\begin{enumerate}[label={[S\arabic*]},leftmargin=*]

\item Emilio Calvano, Giacomo Calzolari, Vincenzo Denicolò, and Sergio Pastorello (2020). \emph{Artificial Intelligence, Algorithmic Pricing, and Collusion}. American Economic Review 110(10), 3267--3297. \url{https://www.aeaweb.org/articles?id=10.1257/aer.20190623}\par \textit{Source role:} Constructive simulation evidence for algorithmic supracompetitive prices. \textit{Locator:} Abstract and article metadata.

\item Stephanie Assad, Robert Clark, Daniel Ershov, and Lei Xu (2024). \emph{Algorithmic Pricing and Competition: Empirical Evidence from the German Retail Gasoline Market}. Journal of Political Economy 132(3), 723--771. \url{https://doi.org/10.1086/726906}\par \textit{Source role:} Market-specific evidence on adoption and competition. \textit{Locator:} Abstract and article metadata.

\item Joseph E. Harrington, Jr. (2018). \emph{Developing Competition Law for Collusion by Autonomous Artificial Agents}. Journal of Competition Law \& Economics 14(3), 331--363. \url{https://doi.org/10.1093/joclec/nhy016}\par \textit{Source role:} Legal-economic conceptual analysis. \textit{Locator:} Abstract and article metadata.

\end{enumerate}

\subsection*{Common conventions: Source mathematical conventions}

Unless an entry states otherwise, agent and firm sets are finite. State and action spaces are measurable, strategies and outcome maps are measurable, probability laws are specified, and expectations exist for the targets under discussion. Infinite-horizon discount factors lie in $(0,1)$ and discounted objectives are finite. Norms, tolerances and complexity input sizes must be specified for computational comparisons. Constants or primitives that appear locally are inputs, not empirical facts asserted by this catalogue.

\textbf{Identification.} A statistical model $m\in\mathcal M$ implies an observable law $P_m$ and target $\theta(m)$. At law $P$, its identified set is
\[
\Theta(P;\mathcal M)=\{\theta(m):m\in\mathcal M,\ P_m=P\}.
\]
At a compatible law $P$, point identification holds when this set is a singleton, equivalently when $\theta(m_1)=\theta(m_2)$ for every $m_1,m_2\in\mathcal M$ with $P_{m_1}=P_{m_2}=P$. Model-wide point identification requires this property at every compatible law, equivalently $P_{m_1}=P_{m_2}$ implies $\theta(m_1)=\theta(m_2)$ for all $m_1,m_2\in\mathcal M$. Sharp bounds mean bounds attained, or approached when the boundary is not attained, by compatible models. Writing this set is a definition, not a solution: the substantive task is to establish informative restrictions and derive or estimate the set. An unrestricted-confounding model may give only trivial bounds.

\textbf{Inference.} Identification, estimability and valid uncertainty quantification are separate questions. Fix $\alpha\in(0,1)$. A confidence procedure $C_n$ over a declared law class $\mathcal P$ has asymptotic uniform coverage at level $1-\alpha$ when
\[
\liminf_{n\to\infty}\inf_{P\in\mathcal P}\Pr_P\{\theta(P)\in C_n\}\geq1-\alpha.
\]
For set-identified targets the coverage event must be defined separately, for example coverage of the full identified set. If an entry uses identification-indexed models instead of uniquely indexed laws, the same coverage convention is applied over those models. Pointwise limits do not establish uniform validity, and asymptotic validity does not guarantee finite-sample precision. Sample sizes, dependence structures and weak-identification sequences are part of each application.

\textbf{Counterfactuals.} $Y_i(a)$ is a potential outcome under a specified intervention, including its timing, implementation and versions. It is not $Y_i$ conditional on voluntarily choosing $a$. A full intervention vector permits interference; shorthand scalar treatment notation does not silently impose no interference. Assignment, selection, attrition, anticipation and measurement must be specified. A claim about an instrument's local effect is not automatically a population policy effect.

\textbf{Economic equilibrium.} A counterfactual model includes best responses, constraints, feasibility, market clearing, state transitions and expectations consistent with the stated information. When several equilibria exist, either a declared selector is an input or the target is set-valued. Comparisons across policy regimes must use the stated selection convention. Reduced-form relationships are not presumed invariant after an equilibrium-changing intervention.

\textbf{Optimization and welfare.} Optimization is over the stated feasible set. If attainment is not established, $\sup$ or $\inf$ and $\varepsilon$-optimal choices replace an asserted maximizer or minimizer. For example, with value $V=\sup_{a\in\mathcal A}W(a)<\infty$, an $\varepsilon$-optimal set is $\{a\in\mathcal A:W(a)\geq V-\varepsilon\}$ for $\varepsilon>0$. Benefits, costs, risk and health or environmental outcomes can be aggregated only using declared units and conversion weights. Interpersonal utility weights, the treatment of fiscal transfers and foreign profits, discounting, and the behavioral welfare criterion are explicit inputs. A comparative-static derivative additionally requires existence and appropriate differentiability of the selected counterfactual.

\textbf{Sources.} Local labels [S1], [S2], and so forth restart in every question. Locators identify the relevant abstract, model, section or result at the level actually checked; a bibliographic or abstract-level verification is not represented as inspection of an entire proof. Working papers are distinguished from later publications and changing versions. Source summaries are paraphrased. Source-role verification does not constitute an exhaustive literature-status audit.
% END ECONOMICS STATEMENT
\end{document}
